\chapter{War conditions: no internet, no GPS, radio only}\label{ch:offline}

In war there is no internet, no GPS and often no phone signal. Near phones sitting near the front are
switched off or in airplane mode, because any emitter can be located. What every unit keeps is the
tactical voice radio. This chapter describes how Near works with nothing else, and how a phone, if one is
available, adds its sensors without a data network.

\section{What runs offline on the laptop}

\begin{table}[htbp]\centering\small
\caption{Every dependency that used to need the internet, and its offline replacement.}\label{tab:offline}
\begin{tabularx}{\linewidth}{>{\raggedright}p{3.4cm}>{\raggedright}p{4.3cm}X}
\toprule
Need & Before (online) & Now (offline)\\
\midrule
Speech-to-text & Chrome recognition (Google servers) & Whisper large-v3-turbo on the GPU, small on CPU; 0.3--0.5 s per transmission\\
Language model & Mistral API & Mistral 7B served locally by Ollama (optional; rules always run)\\
Map background & tile.openstreetmap.org & 682 tiles drawn on the laptop from our own OSM and IGN data (13 s, zooms 13--17)\\
Map and 3D libraries & unpkg CDN & Bundled in \code{frontend/vendor/}\\
Sound models & Hugging Face downloads & Cached once; loaded with \code{local\_files\_only}\\
Spoken replies & Browser voice, default speaker only & Windows SAPI / espeak-ng to WAV, played on the radio's audio output\\
\bottomrule
\end{tabularx}
\end{table}

One command, run while there is still internet, downloads everything (\code{scripts/fetch\_models.py}); a
second draws the map tiles (\code{scripts/render\_tiles.py}). Rendering our own tiles from data we already
hold also respects OpenStreetMap's tile policy, which forbids bulk downloads.

\section{Radio mode}

\nanofig{radiochain}{Radio mode: the laptop listens to the radio, understands each transmission, and answers on it after the
operator approves.}
{Horizontal signal-chain diagram, 16:9, plain white background, flat vector, navy outlines, amber audio waves, cyan data arrows.
Left: a soldier with a handheld radio, a speech bubble ``This is Bravo 2...''. Zig-zag radio waves to a second radio on a
table at the CCP. An audio cable from the radio's speaker to a laptop labelled ``audio in''. Inside a big laptop outline, a row
of boxes: ``Voice activity: cut transmission'' → ``Data tones? decode packet'' → ``Whisper (local)'' → ``Callsign: BRAVO 2''
→ ``Near pipeline'' → ``Reply draft''. Below the row, a small operator icon with a check mark labelled ``approve''. Then
``Offline voice (SAPI)'' → an audio cable labelled ``audio out'' back to the radio's microphone, and radio waves back to the
soldier with a bubble ``Bravo 2, this is CCP. Copy... Over.''. A crossed-out cloud icon labelled ``no internet''. Exact labels. Colour palette (use only these): navy \#0D1110, cyan \#1AA7C4, amber \#E09A2B, red \#E5484D, green \#3FA877, grey \#8B9A94, plus white for text on dark shapes.}

\paragraph{Wiring.} The radio's speaker output goes into the laptop's audio input, and the laptop's audio output into the
radio's microphone. In the browser, the operator picks both devices on the Radio panel. Replies are played with
\code{setSinkId} on the chosen output, so the laptop's own speaker stays free.

\paragraph{Cutting transmissions.} The page tracks the hiss level when nobody talks (slow average of the RMS). A
transmission starts when the level exceeds four times that floor (at least 0.01). It ends after 900 ms of quiet, or after
30 s. Clips shorter than 0.5 s are dropped. While the laptop is transmitting a reply, listening pauses so it doesn't hear
itself.

\paragraph{Understanding one transmission} (\code{radio.py}):
\begin{enumerate}
  \item \textbf{Data tones} are decoded first (next section). A measurement packet becomes the same evidence as a phone's Listen
        clip.
  \item \textbf{Speech} is transcribed by local Whisper. Transcripts that are only noise words are dropped.
  \item \textbf{Callsign.} ``This is Bravo two'', ``ici Charlie 3'', or a bare ``Bravo 2. Left. Over.'' gives \code{BRAVO 2}.
        A known callsign goes to its open report; a new one opens a report. A transmission with no callsign goes to the
        report heard on the radio in the last 15 minutes. Procedure words (``this is'', ``over'', ``out'', ``roger'') are
        stripped before parsing.
  \item \textbf{Threats in the background sound.} The radio carries more than the voice. Gunfire and explosions come from the
        big laptop models; drones from Near-small, which was trained on radio audio (the big drone model fails there:
        Table~\ref{tab:radioeval}). They're flagged ``heard over the radio''.
  \item \textbf{Reply.} ``Bravo 2, this is CCP.'' + the read-back + ``Over.'' (French: \emph{ici poste de secours} ...
        \emph{À vous}). The operator can edit, send or skip. ``Send without asking'' makes it automatic.
\end{enumerate}

\begin{keyresult}[End-to-end test through a simulated analog radio]
A full spoken report was transcribed word for word, including ``tourniquet on at 1432'' and ``a church at bearing 240
degrees'', with callsign \code{BRAVO 2}. A 25-byte packet sent as tones in the same session added a compass heading, the
kind of space, two wall echoes, the noise level, traffic and a drone threat to the same casualty. A drone recording sent over
the radio raised a drone alert (0.99) on the active casualty. A short answer, ``Bravo 2. Left. Over.'', was routed to the
right report.
\end{keyresult}

\section{Data over the voice channel: the modem}

A phone in airplane mode can still make sound. \code{modem.py} and its JavaScript twin \code{modem.js} turn a few bytes
into audio tones that any voice radio carries. Amateur radio's APRS has sent data this way over FM voice radios for decades.

\begin{table}[htbp]\centering\small
\caption{Modem parameters.}\label{tab:modem}
\begin{tabular}{ll}
\toprule
Modulation & binary FSK, continuous phase\\
Tones & 1200 Hz = 1, 2200 Hz = 0 (inside the 300--3400 Hz radio band)\\
Rate & 400 bit/s (integer samples per bit at 8, 16 and 48 kHz)\\
Frame & preamble 0x55 $\times$ 6 $|$ then twice: sync 0x2D 0xD4 $|$ length $|$ payload $|$ CRC-16 (CCITT)\\
Packet of 40 bytes & $\approx$ 2.0 s of tones with two copies (1.1 s with one)\\
Decode time & 3 ms for a 2 s packet (Python)\\
\bottomrule
\end{tabular}
\end{table}

\paragraph{Decoding.} For each tone, the receiver mixes the signal with $e^{-j2\pi f t}$ and takes the magnitude of its sum
over one bit (a sliding Goertzel-style correlator, computed with a cumulative sum). The soft bit is the mark energy minus the
space energy, normalised. Because the receiver doesn't know where a bit starts, it tries 8 phase offsets per bit. It looks
for the 16-bit sync word with at most one bit wrong, and reads the length byte. The length is also inferred from the gap
between two sync words, in case the byte itself was hit. Each candidate payload is checked against its CRC.

\paragraph{Two copies, added up.} Digital radios lose short frames of audio (about 20 ms each) on a bad link. A lost frame
wipes out 8 bits, so a single-copy packet of 1.1 s rarely survives. Each frame is therefore sent twice, and the receiver adds
the two copies' soft bits before deciding. A dropout rarely hits the same bits in both. This is soft-decision diversity
combining. It's simpler than Reed--Solomon coding, and it's easy to mirror byte for byte in the browser.

\begin{figure}[htbp]\centering
\begin{tikzpicture}
\begin{axis}[ybar, width=0.95\linewidth, height=6cm, bar width=9pt, ymin=0, ymax=45,
  ylabel={Packets decoded (of 40)}, symbolic x coords={clean,analog 20 dB,analog 10 dB,digital 20 dB,digital 12 dB},
  xtick=data, x tick label style={font=\footnotesize}, legend style={at={(0.5,1.03)},anchor=south,legend columns=-1},
  enlarge x limits=0.12]
\addplot[fill=greyN!60] coordinates {(clean,40) (analog 20 dB,40) (analog 10 dB,40) (digital 20 dB,10) (digital 12 dB,8)};
\addplot[fill=cyanN] coordinates {(clean,40) (analog 20 dB,40) (analog 10 dB,40) (digital 20 dB,39) (digital 12 dB,39)};
\addplot[fill=amberN] coordinates {(clean,40) (analog 20 dB,40) (analog 10 dB,40) (digital 20 dB,40) (digital 12 dB,39)};
\legend{1 copy (1.1 s), 2 copies (2.0 s), 3 copies (2.9 s)}
\end{axis}
\end{tikzpicture}
\caption{Random 40-byte packets through the simulated radio channel, random start time and level.}\label{fig:modemchart}
\end{figure}

\nanofig{modem}{The modem: bytes become two audio tones; the frame is sent twice and the two copies are added before deciding
each bit.}
{Technical signal illustration, 16:9, plain white background, flat vector. Top: a row of byte boxes ``55 55 55 55 55 55'' (grey,
labelled ``preamble''), then ``2D D4'' (navy, ``sync''), ``1C'' (``length''), a long cyan block ``payload 28 B'', ``CRC'' (amber),
then the same sync-length-payload-CRC sequence again labelled ``copy 2''. Middle: an audio waveform below the frame alternating
between a lower-frequency sine labelled ``1200 Hz = 1'' and a higher-frequency sine labelled ``2200 Hz = 0''. Bottom: two
rows of soft-bit bars labelled ``copy 1'' and ``copy 2'', each with a short gap labelled ``20 ms dropout'' at different places,
plus a third row labelled ``sum'' with no gap. Exactly these labels. Colour palette (use only these): navy \#0D1110, cyan \#1AA7C4, amber \#E09A2B, red \#E5484D, green \#3FA877, grey \#8B9A94, plus white for text on dark shapes.}

\section{Compact packets}

\code{packet.py} and \code{packet.js} share one byte layout. A packet carries what the phone measured, never raw audio.

\begin{table}[htbp]\centering\small
\caption{Report packet (phone $\rightarrow$ CCP). Fields after the mask are present only if their bit is set.}\label{tab:packet}
\begin{tabularx}{\linewidth}{llX}
\toprule
Bit & Field & Encoding\\
\midrule
-- & type / version & 1 B: 0x11\\
-- & callsign & 1 B length + up to 8 ASCII bytes\\
-- & time and sequence & 2 B: minutes since midnight UTC (11 bits) + sequence (5 bits)\\
-- & field mask & 2 B\\
0 & compass (``12 o'clock'') & 1 B: degrees $\times$ 256 / 360\\
1 & aimed landmarks & 1 B count + 2 B each (azimuth, elevation $\times$ 4)\\
2 & step count & 2 B\\
3 & kind of space & 2 B: 4 probabilities $\times$ 15, 4 bits each\\
4 & noise level & 1 B: dB(A)\\
5 & wall echoes & 1 B count + 1 B each (metres $\times$ 4, 0.25 m steps up to 63.75 m)\\
6 & sounds heard & 1 B count + 1 B each (group 3 bits, score 5 bits)\\
7 & threats heard & 1 B mask (drone, small arms, explosion) + 1 B score each\\
8 & guidance event & 2 B: reached / arrived / lost + leg number\\
9 & strength of the space clues & 1 B\\
\bottomrule
\end{tabularx}
\end{table}

A full measurement packet is 21--35 bytes; the example in Chapter~\ref{ch:eval} carries all ten fields in 32 bytes. The route
packet (CCP $\rightarrow$ phone) is 0x21, callsign, route id, leg count, then 6 bytes per leg: bearing (1 B), length (2 B),
landmark index (2 B, from a list the phone keeps offline) and the side of that landmark (1 B). An 8-leg route is 53 bytes,
about 2.5 s of tones. The same bytes can travel over a radio's data mode or a mesh radio, posted to \code{/api/packet}.

\nanofig{packet}{A 32-byte measurement packet, field by field.}
{Byte-map infographic, 16:9, plain white background, flat vector. A long horizontal strip of 32 small square cells grouped and
coloured by field, with labels above each group: ``type'' (1), ``callsign BRAVO 2'' (8), ``time+seq'' (2), ``mask'' (2),
``compass 75°'' (1), ``aim'' (3), ``steps 1234'' (2), ``space'' (2), ``dB(A) 62'' (1), ``walls 3.5 m, 9 m'' (3), ``traffic''
(2), ``drone 0.85'' (2), ``reached leg 2'' (2), ``strength'' (1). Below the strip, a bracket labelled ``32 bytes = about 2 s of
tones''. Pastel colours per group, navy labels. Exactly these labels. Colour palette (use only these): navy \#0D1110, cyan \#1AA7C4, amber \#E09A2B, red \#E5484D, green \#3FA877, grey \#8B9A94, plus white for text on dark shapes.}

\section{The radio channel we simulate}

\code{radiosim.py} approximates what a radio does to sound, for testing and for training:
\begin{itemize}
  \item \textbf{Analog FM:} band-pass 300--3400 Hz (4th-order Butterworth), automatic gain to a fixed level, soft clipping
        ($\tanh$), resampling to 8 kHz, hiss at 15--30 dB SNR.
  \item \textbf{Digital:} the same, plus 8-bit $\mu$-law quantisation and 3 \% of 20 ms frames lost.
\end{itemize}
Real military radios use speech vocoders (e.g. MELPe) at low bit rates, which model speech and may smear tones and other
sounds further. That's why every radio result in this report needs to be repeated on the unit's actual radio.

\nanofig{radiosim}{The simulated radio channel used for tests and training.}
{Signal-processing block diagram, 16:9, plain white background, flat vector. Left: a clean colourful spectrogram labelled ``phone
microphone''. Then boxes in a row connected by arrows: ``band-pass 300–3400 Hz'', ``automatic gain + soft clipping'',
``8 kHz'', ``8-bit µ-law (digital)'', ``lost 20 ms frames (digital)'', ``hiss 15–30 dB''. Right: the same spectrogram now narrow
in frequency, with two black vertical gaps and grainy noise, labelled ``what the CCP hears''. Exactly these labels. Colour palette (use only these): navy \#0D1110, cyan \#1AA7C4, amber \#E09A2B, red \#E5484D, green \#3FA877, grey \#8B9A94, plus white for text on dark shapes.}

\section{The phone with no network}

If the unit has a phone, it adds sensors the radio can't carry. The field page installs as an offline web app: a service
worker caches the page, the model, the libraries and the landmark list on first load. We tested it by reloading with the
network cut.

\nanofig{phone}{The offline phone: measure, send as tones, receive a route, follow the arrow.}
{Smartphone UI mockup, portrait, dark theme, shown on a plain white background, 9:16 phone inside a 16:9 canvas with
callouts. Screen content top to bottom: header ``NEAR · FIELD'', a card titled ``No network: talk on the radio, send data as
tones'', a text field ``Bravo 2'', a button ``Measure 4 s'', result chips ``outdoors, between buildings 72\%'', ``walls 3.5, 9 m'',
``drone 0.91'', ``62 dB(A)'', a large button ``Send by radio'', two buttons ``Receive route'' and ``Hand-over QR''. Below, a
guidance panel with a big cyan arrow pointing up-right, a chip ``86 m to the next turn'', text ``Leg 2 of 5: bearing 043, about 120
metres. You should see Mairie de Saint-Mandé on your right.'' and a button ``Reached the turn''. Callouts outside the phone:
``runs offline'', ``~30 bytes as 2 s of tones'', ``compass + step counter''. Exactly these labels. Colour palette (use only these): navy \#0D1110, cyan \#1AA7C4, amber \#E09A2B, red \#E5484D, green \#3FA877, grey \#8B9A94, plus white for text on dark shapes.}

\begin{itemize}
  \item \textbf{Measure 4 s.} The phone records 3.1 s of ambient sound, then plays the 0.5 s chirp. Near-small runs in the browser
        (onnxruntime-web, WebAssembly, one thread) and returns threats, sounds and the kind of space. A matched filter on the chirp
        (FFT cross-correlation with the same sweep the phone played) finds the direct sound and up to four wall echoes between
        1.8 and 25 m. The noise level comes from the quiet part, calibrated on the operator's own voice from the last
        push-to-talk (the loudest 20 \% of 50 ms frames taken as about 62 dB(A)).
  \item \textbf{Send by radio.} The operator presses push-to-talk, holds the phone's speaker to the radio's microphone, and
        taps: the packet plays as tones.
  \item \textbf{Receive route.} The phone listens to the radio's speaker, decodes tones every 2.5 s, and starts guiding when
        a route packet arrives. An arrow points along the leg's bearing relative to the compass heading. The step counter
        (accelerometer peaks above the running mean, at least 0.3 s apart, 0.7 m per step) counts down the metres. The turns
        are spoken, with landmark cues from the cached list. ``Reached'' advances the leg and queues a guidance event for
        the next packet.
  \item \textbf{Hand-over QR.} Everything the phone knows (callsign, time, last read-back, threats, facing) as a QR code for the
        medic who arrives.
\end{itemize}

The phone and the laptop were checked against each other in both directions. A packet encoded in JavaScript was decoded
in Python after the simulated digital radio. A route encoded in Python was decoded in JavaScript after the simulated
analog radio.

\section{Soldiers without a phone}

Most of what a phone adds has a spoken equivalent, so a soldier with only a radio loses little:
\begin{itemize}
  \item \textbf{Compass and Aim \& tap $\rightarrow$ spoken compass bearings.} Every soldier carries a compass, and resection
        from bearings is standard land navigation. Two named landmarks give a 71 m median error (Chapter~\ref{ch:loc}).
  \item \textbf{Step counter $\rightarrow$ pace count.} Soldiers are trained to count paces per 100 m; ``300 metres done''
        moves the pin.
  \item \textbf{Phone sound analysis $\rightarrow$ the radio audio itself.} Threats are detected at the CCP in what comes
        through the radio.
  \item \textbf{Echo chirp:} lost. A half-duplex radio can't send a chirp and listen to its echo. The question
        ``indoors, in a street, or in the open?'' partly replaces it.
\end{itemize}
